% GENERATED FILE. Edit canonical lessons, then run npm run generate:books.
% canonical-source-hash: fnv1a64:ac6c3b0e5652a3fd
% canonical-lessons: HI-C289-jebkatra, HI-C289-chinna, HI-C289-jhapatna, HI-C289-farar, HI-C289-picha

\chapter{The Thief Got Away}
\label{ch:hi-the-thief-got-away}
\hlchaptermodality{289}

\begin{chapteropening}
\textbf{By the end of this chapter:} \emph{I can say how a thief got away and who went after him.}

Complete the last lesson of chapter 289: I can say how a thief got away and who went after him.
\end{chapteropening}

\section[\texorpdfstring{\hi{जेबकतरा} (jebkatrā)}{जेबकतरा (jebkatrā)}]{\hi{जेबकतरा} (jebkatrā) — a pickpocket}

\label{lesson:HI-C289-jebkatra}



\begin{quote}
Before the new one: say the Hindi for a witness, then the Hindi for outside.
\end{quote}

\subsection*{You'll want to know: \hi{जेबकतरा}}
\textbf{\hi{जेबकतरा}} — \emph{jebkatrā} — \textquotedblleft{}a pickpocket\textquotedblright{}.

\begin{quote}
\hi{बाज़ार} \hi{में} \hi{एक} \hi{जेबकतरा} \hi{घूम} \hi{रहा} \hi{था।}
\end{quote}

\begin{quote}
\emph{bāzār meṁ ek jebkatrā ghūm rahā thā.}
\end{quote}

\begin{quote}
A pickpocket was wandering through the market.
\end{quote}

Masculine.

\begin{cousinweb}
\textbf{\hi{जेबकतरा}} (\emph{jebkatrā}) is \textbf{\hi{जेब}}, a pocket, and \emph{katrā}, one who snips, from \emph{katarnā}, to clip, which CDIAL 2858 lists for Hindi under Sanskrit \emph{kartarī}, \textquotedblleft{}scissors\textquotedblright{}: the one who cuts into pockets.
\end{cousinweb}

\begin{grammarlens}[title={What you've built}]
One more word for telling what happened.
\end{grammarlens}

\subsection*{Your turn}
\begin{itemize}
\raggedright
  \item \emph{Say it:} \emph{jebkatrā}
  \item \emph{Say it:} \emph{jebkatrā}, once more
  \item \emph{Say it:} \emph{bāzār meṁ ek jebkatrā ghūm rahā thā.} — A pickpocket was wandering through the market.
  \item \emph{Recall:} say the Hindi for a witness, then the Hindi for outside, then say \emph{jebkatrā} again
\end{itemize}

\begin{tcolorbox}[breakable,colback=teal!4,colframe=teal!35!black,title={Before you move on}]
What does \hi{जेबकतरा} mean? (\textquotedblleft{}A pickpocket\textquotedblright{}.) Say it once more.
\end{tcolorbox}
\section[\texorpdfstring{\hi{छीनना} (chīnnā)}{छीनना (chīnnā)}]{\hi{छीनना} (chīnnā) — to snatch}

\label{lesson:HI-C289-chinna}



\begin{quote}
Before the new one: say the Hindi for outside, then the Hindi for a pickpocket.
\end{quote}

\subsection*{You'll want to know: \hi{छीनना}}
\textbf{\hi{छीनना}} — \emph{chīnnā} — \textquotedblleft{}to snatch\textquotedblright{}.

\begin{quote}
\hi{उसने} \hi{मीरा} \hi{के} \hi{हाथ} \hi{से} \hi{मोबाइल} \hi{छीन} \hi{लिया।}
\end{quote}

\begin{quote}
\emph{usne Mīrā ke hāth se mobāil chīn liyā.}
\end{quote}

\begin{quote}
He snatched the mobile phone out of Mira's hand.
\end{quote}

Past, with \hi{ने} (he or she …): \hi{उसने} \hi{थैला} \hi{छीना}, \emph{usne thailā chīnā}.

\begin{cousinweb}
\textbf{\hi{छीनना}} (\emph{chīnnā}) is set by CDIAL under Sanskrit \emph{chinna}, \textquotedblleft{}cut off\textquotedblright{} (CDIAL 5047): what is snatched is torn away from its owner. It often takes \textbf{\hi{लेना}}: \textbf{\hi{छीन} \hi{लिया}}, snatched it away.
\end{cousinweb}

\begin{grammarlens}[title={What you've built}]
One more word for telling what happened.
\end{grammarlens}

\subsection*{Your turn}
\begin{itemize}
\raggedright
  \item \emph{Say it:} \emph{chīnnā}
  \item \emph{Say it:} \emph{chīnnā}, once more
  \item \emph{Say it:} \emph{usne Mīrā ke hāth se mobāil chīn liyā.} — He snatched the mobile phone out of Mira's hand.
  \item \emph{Recall:} say the Hindi for outside, then the Hindi for a pickpocket, then say \emph{chīnnā} again
\end{itemize}

\begin{tcolorbox}[breakable,colback=teal!4,colframe=teal!35!black,title={Before you move on}]
What does \hi{छीनना} mean? (\textquotedblleft{}To snatch\textquotedblright{}.) Say it once more.
\end{tcolorbox}
\section[\texorpdfstring{\hi{झपटना} (jhapaṭnā)}{झपटना (jhapaṭnā)}]{\hi{झपटना} (jhapaṭnā) — to lunge; to swoop}

\label{lesson:HI-C289-jhapatna}



\begin{quote}
Before the new one: say the Hindi for a pickpocket, then the Hindi for to snatch.
\end{quote}

\subsection*{You'll want to know: \hi{झपटना}}
\textbf{\hi{झपटना}} — \emph{jhapaṭnā} — \textquotedblleft{}to lunge; to swoop\textquotedblright{}.

\begin{quote}
\hi{चोर} \hi{मेरे} \hi{थैले} \hi{पर} \hi{झपटा} \hi{और} \hi{भाग} \hi{गया।}
\end{quote}

\begin{quote}
\emph{cor mere thaile par jhapṭā aur bhāg gayā.}
\end{quote}

\begin{quote}
The thief lunged at my bag and ran off.
\end{quote}

Past (he …): \hi{वह} \hi{थैले} \hi{पर} \hi{झपटा}, \emph{vah thaile par jhapṭā}.

\begin{cousinweb}
\textbf{\hi{झपटना}} (\emph{jhapaṭnā}) comes from \emph{*jhapp}, a sudden movement, a form CDIAL reconstructs (CDIAL 5336). A hawk \textbf{\hi{झपटता} \hi{है}} on its prey; a thief \textbf{\hi{झपटता} \hi{है}} on a bag: \textbf{\hi{थैले} \hi{पर} \hi{झपटा}}.
\end{cousinweb}

\begin{grammarlens}[title={What you've built}]
One more word for telling what happened.
\end{grammarlens}

\subsection*{Your turn}
\begin{itemize}
\raggedright
  \item \emph{Say it:} \emph{jhapaṭnā}
  \item \emph{Say it:} \emph{jhapaṭnā}, once more
  \item \emph{Say it:} \emph{cor mere thaile par jhapṭā aur bhāg gayā.} — The thief lunged at my bag and ran off.
  \item \emph{Recall:} say the Hindi for a pickpocket, then the Hindi for to snatch, then say \emph{jhapaṭnā} again
\end{itemize}

\begin{tcolorbox}[breakable,colback=teal!4,colframe=teal!35!black,title={Before you move on}]
What does \hi{झपटना} mean? (\textquotedblleft{}To lunge; to swoop\textquotedblright{}.) Say it once more.
\end{tcolorbox}
\section[\texorpdfstring{\hi{फ़रार} (farār)}{फ़रार (farār)}]{\hi{फ़रार} (farār) — on the run}

\label{lesson:HI-C289-farar}



\begin{quote}
Before the new one: say the Hindi for to snatch, then the Hindi for to lunge.
\end{quote}

\subsection*{You'll want to know: \hi{फ़रार}}
\textbf{\hi{फ़रार}} — \emph{farār} — \textquotedblleft{}on the run\textquotedblright{}.

\begin{quote}
\hi{चोर} \hi{दो} \hi{हफ़्ते} \hi{से} \hi{फ़रार} \hi{है।}
\end{quote}

\begin{quote}
\emph{cor do hafte se farār hai.}
\end{quote}

\begin{quote}
The thief has been on the run for two weeks.
\end{quote}

\begin{cousinweb}
\textbf{\hi{फ़रार}} (\emph{farār}) is usually traced, through Persian, to Arabic \emph{firār}, \textquotedblleft{}flight\textquotedblright{}. It describes the person, not the running: \textbf{\hi{चोर} \hi{फ़रार} \hi{है}}, the thief got away and has not been found.
\end{cousinweb}

\begin{grammarlens}[title={What you've built}]
One more word for telling what happened.
\end{grammarlens}

\subsection*{Your turn}
\begin{itemize}
\raggedright
  \item \emph{Say it:} \emph{farār}
  \item \emph{Say it:} \emph{farār}, once more
  \item \emph{Say it:} \emph{cor do hafte se farār hai.} — The thief has been on the run for two weeks.
  \item \emph{Recall:} say the Hindi for to snatch, then the Hindi for to lunge, then say \emph{farār} again
\end{itemize}

\begin{tcolorbox}[breakable,colback=teal!4,colframe=teal!35!black,title={Before you move on}]
What does \hi{फ़रार} mean? (\textquotedblleft{}On the run\textquotedblright{}.) Say it once more.
\end{tcolorbox}
\section[\texorpdfstring{\hi{पीछा} (pīchā)}{पीछा (pīchā)}]{\hi{पीछा} (pīchā) — a chase; the pursuit}

\label{lesson:HI-C289-picha}



\begin{quote}
Before the new one: say the Hindi for to lunge, then the Hindi for on the run.
\end{quote}

\subsection*{You'll want to know: \hi{पीछा}}
\textbf{\hi{पीछा}} — \emph{pīchā} — \textquotedblleft{}a chase; the pursuit\textquotedblright{}.

\begin{quote}
\hi{अगर} \hi{मेरे} \hi{पास} \hi{साइकिल} \hi{होती} \hi{तो} \hi{मैं} \hi{चोर} \hi{का} \hi{पीछा} \hi{करता।}
\end{quote}

\begin{quote}
\emph{agar mere pās sāikil hotī to maiṁ cor kā pīchā kartā.}
\end{quote}

\begin{quote}
If I had had a bicycle, I would have gone after the thief.
\end{quote}

Masculine.

\begin{cousinweb}
\textbf{\hi{पीछा}} (\emph{pīchā}) is the noun behind \textbf{\hi{पीछे}}, behind, from \emph{*paśca}, the back part, a form CDIAL reconstructs (CDIAL 7990). \textbf{\hi{किसी} \hi{का} \hi{पीछा} \hi{करना}} is to go after someone: to make yourself the one behind them.
\end{cousinweb}

\begin{grammarlens}[title={What you've built}]
One more word for telling what happened.
\end{grammarlens}

\subsection*{Your turn}
\begin{itemize}
\raggedright
  \item \emph{Say it:} \emph{pīchā}
  \item \emph{Say it:} \emph{pīchā}, once more
  \item \emph{Say it:} \emph{agar mere pās sāikil hotī to maiṁ cor kā pīchā kartā.} — If I had had a bicycle, I would have gone after the thief.
  \item \emph{Recall:} say the Hindi for to lunge, then the Hindi for on the run, then say \emph{pīchā} again
\end{itemize}

\begin{tcolorbox}[breakable,colback=teal!4,colframe=teal!35!black,title={Before you move on}]
What does \hi{पीछा} mean? (\textquotedblleft{}A chase; the pursuit\textquotedblright{}.) Say it once more.
\end{tcolorbox}
